\documentclass{article}

\usepackage[margin=0.36in]{geometry}
\usepackage[T1]{fontenc}
\usepackage[utf8]{inputenc}
\usepackage{helvet}
\usepackage{xcolor}
\usepackage{tabularx}
\usepackage{booktabs}
\usepackage{array}
\usepackage{enumitem}
\usepackage[hidelinks]{hyperref}

\renewcommand{\familydefault}{\sfdefault}
\pagestyle{empty}
\setlength{\parindent}{0pt}
\setlength{\parskip}{0.8pt}
\setlength{\tabcolsep}{3.6pt}
\setlist[itemize]{leftmargin=*, nosep, topsep=1pt}
\renewcommand{\arraystretch}{1.0}

\definecolor{navy}{HTML}{081B33}
\definecolor{blue}{HTML}{0F67B1}
\definecolor{green}{HTML}{14784D}
\definecolor{ink}{HTML}{1F2933}
\definecolor{muted}{HTML}{5D6B78}

\newcolumntype{Y}{>{\raggedright\arraybackslash}X}

\newcommand{\labeltext}[1]{{\scriptsize\bfseries\color{muted}\MakeUppercase{#1}}}
\newcommand{\sectionhead}[1]{\vspace{1.4pt}{\color{navy}\bfseries\MakeUppercase{#1}}\par{\color{blue}\rule{\linewidth}{0.4pt}}\vspace{0.3pt}}
\newcommand{\metricline}[3]{%
  \begin{minipage}[t]{0.315\linewidth}
    \labeltext{#1}\par
    {\color{#2}\fontsize{16}{17}\selectfont\bfseries #3}
  \end{minipage}
}

\begin{document}
\enlargethispage{2\baselineskip}
\fontsize{7.85}{8.85}\selectfont

\begin{minipage}[t]{\linewidth}
{\color{navy}\fontsize{22}{23}\selectfont\bfseries NUWA Evidence Sheet}\par
{\color{ink}\normalsize Autonomous self-healing ecosystem for low Earth orbit satellite constellations}\par
{\color{muted}\small Welcome to NUWA, the first complete, autonomous, self-healing ecosystem for satellite constellations. NUWA transforms a group of fragile, isolated spacecraft into a cooperative, resilient network.}
\end{minipage}

\vspace{4pt}
{\color{blue}\rule{\linewidth}{0.7pt}}
\vspace{1pt}

\metricline{Primary technical KPI}{blue}{94.3\%}
\hfill
\metricline{Primary sustainability KPI}{green}{3,103 t}
\hfill
\metricline{Protocol test pass rate}{blue}{273/273}

\vspace{-1pt}
\begin{tabularx}{\linewidth}{@{}Y Y Y@{}}
{\scriptsize RMSE improvement in IT-05 30-day drift: raw 8.909 to corrected 0.504.} &
{\scriptsize CO$_2$-eq avoided per year in the 100-satellite reference model.} &
{\scriptsize C++ tests pass across codec, frame pipeline, FSM, DEGR, simulation, and matrix tests.}
\end{tabularx}

\vspace{-1pt}
\begin{tabularx}{\linewidth}{@{}Y@{\hspace{10pt}}Y@{}}
\sectionhead{Problem Statement}
Low Earth orbit is crowded, fragile, and time-critical. A sensor can drift, telemetry can become unreliable, and a satellite can lose confidence in its own data before ground teams are available. Recoverable faults then become early decommission decisions, adding replacement pressure, launch emissions, and long-term debris exposure. NUWA targets that gap: satellites are not always dead; sometimes they are only blind.
&
\sectionhead{Core Function}
NUWA gives satellites neighbors. It runs on the edge as an autonomous self-healing ecosystem: Hankel-SVD anomaly detection catches sensor faults, then the SISP protocol coordinates correction, relay, and borrow services through fixed 64-byte frames and a deterministic 21-state/24-event C++ state machine. A degraded satellite can correct drift, relay vital data through a peer, or borrow a healthy neighbor sensor stream.
\end{tabularx}

\sectionhead{Baseline Case vs. Improved Case}
\begin{tabularx}{\linewidth}{@{}>{\bfseries\raggedright\arraybackslash}p{0.14\linewidth}Y Y@{}}
\toprule
Impact area & Baseline case & Improved case with NUWA \\
\midrule
Sensor faults & Degraded satellites wait for ground intervention or onboard redundancy. & Neighbor-assisted correction preserves useful measurements during drift or failure. \\
Data delivery & Missed ground contact delays payloads and increases mission latency. & Relay moves data through a better-positioned satellite when direct contact is unavailable. \\
Energy use & Resilience often depends on extra hardware, repeated contacts, or delayed recovery cycles. & Full daily protocol overhead is about 317 mWh/day, or 0.26\% of a 5 W continuous budget. \\
Satellite life cycle & Early sensor or subsystem degradation can shorten usable mission life. & Reference model assumes 45\% life extension, improving useful life from 3.0 to 4.35 years. \\
Replacement pressure & Premature failure increases rebuild, launch, and operating cost. & Replacement launches fall from 33.3/yr to 23.0/yr in the 100-satellite model, about 31\% fewer. \\
Debris over years & More frequent replacements increase mass sent to orbit and long-term debris risk. & About 10.3 future debris objects avoided per year; about 27,800 over 50 years under growth assumptions. \\
Failure behavior & Fault messages can cause poor distributed logic to spread failure state. & Healthy nodes record failed peers without cascading into critical failure. \\
\bottomrule
\end{tabularx}

\vspace{1pt}
\begin{tabularx}{\linewidth}{@{}Y@{\hspace{10pt}}Y@{}}
\sectionhead{Test Conditions}
\begin{itemize}
  \item IT-05: 30-day drift correction using Kalman filter.
  \item IT-06: 7-day scenario with 10\% packet loss and 5 satellites.
  \item OPS-SAT stream test: 4 channels mapped to LEO-IMG, LEO-COM, LEO-SCI, and LEO-OBS.
  \item Relay resilience: checksum drop, out-of-order fragments, retry, and duplicate replay.
  \item Dual-PHY: 437 MHz control at 9,600 bps and bulk at 19,200 bps.
\end{itemize}
&
\sectionhead{Measured Evidence}
\begin{itemize}
  \item IT-05: RMSE reduced from 8.909 to 0.504, a 94.3\% improvement.
  \item IT-06: RMSE reduced from 8.290 to 1.197, an 85.6\% improvement.
  \item 1 MiB bulk relay completes in about 7.91 min at 1.65 Wh link energy.
  \item Dual-PHY validation passed 8/8 assertions.
  \item 8-neighbor correction round takes about 1.10 s and 3.66 J requester energy.
\end{itemize}
\end{tabularx}

\vspace{-1pt}
\begin{tabularx}{\linewidth}{@{}Y@{\hspace{10pt}}Y@{}}
\sectionhead{Technical Basis}
\begin{itemize}
  \item Fixed 512-bit frame, CRC-8/MAXIM checksum, typed payloads, bounded state transitions.
  \item DEGR trust weight: $w_i=\max(0.05,1-\mathrm{DEGR}_i/15)$.
  \item Weighted average, weighted median, Kalman, and hybrid correction filters.
  \item Hankel-SVD stream detector plus OPS-SAT static SVD study reaching up to 0.99 ROC-AUC.
  \item 437 MHz link budget: 8.9 dB margin at 1000 km; about 2,790 km usable range.
\end{itemize}
&
\sectionhead{Project Limitations}
\begin{itemize}
  \item Research-grade implementation, not yet flight-qualified software.
  \item Sustainability outcomes are model estimates and must be tuned per mission.
  \item RF performance depends on geometry, antennas, spectrum authorization, Doppler, and radio hardware.
  \item Security hardening still needs authenticated frames and signed health beacons.
  \item Deorbit coordination is conceptual; next step is COTS-radio hardware-in-the-loop testing.
\end{itemize}
\end{tabularx}

\sectionhead{Primary SDG Alignment}
\begin{tabularx}{\linewidth}{@{}Y Y Y@{}}
{\scriptsize\textbf{SDG 9.} Self-healing satellite infrastructure; fault detection under 1 min and correction under 5 s.} &
{\scriptsize\textbf{SDG 13.} Protects Earth-observation data integrity; telemetry validation above 99\% on nominal satellites.} &
{\scriptsize\textbf{SDG 17.} Reliable space infrastructure for global development data; target above 90\% autonomous fault resolution.}
\end{tabularx}

\sectionhead{Scalability and Future Work}
\begin{tabularx}{\linewidth}{@{}Y Y@{}}
{\scriptsize\textbf{Scale capabilities:} shared computation; low-latency fault detection; multi-frequency state machine; encryption layer.} &
{\scriptsize\textbf{Roadmap:} 2026 validation through multi-topology stress testing; 2027 scale through hardware-in-loop testing and inter-agency sharing protocol; 2028+ deployment and enhancements.}
\end{tabularx}

\sectionhead{Impact Summary}
NUWA keeps degraded satellites useful, holds daily protocol energy to about 0.26\% of a 5 W budget, extends useful satellite life, and reduces long-term replacement, launch-mass, and debris pressure.

{\scriptsize\color{muted}\textbf{Evidence base:} NUWA docs, SDG impact metrics, test summaries, KPI snapshot, signal/energy studies, and equation reference. Measured results are test-derived; sustainability results are scenario estimates.}

\end{document}
